\section{Introduction}\label{sec:intro}
 
Large Language Models (LLMs) are increasingly used as a source of information, including for current events and their interpretation, especially among young users \citep{reuters2026dnr, niemanlab2026chatbots}. Although language models typically have a training data cutoff date and may not be up to date with the latest developments in the world \citep{vu2024freshllms}, web search and retrieval-augmented generation (RAG) have greatly improved their ability to answer queries about recent events and provide factually grounded information \citep{lewis2020rag}. Nevertheless, LLMs remain susceptible to hallucinations \citep{huang2025hallucination}, and research shows that the news integrity of LLM agents can be seriously compromised by sourcing problems and accuracy issues \citep{ebu2025newsintegrity, jazwinska2025citation}. In addition, compared to traditional search engines, where links are explicitly displayed and used more often \citep{reuters2026chatbots}, sources in LLM-synthesized information are often opaque or absent \citep{liu2023verifiability}, and users are less likely to visit linked sources even when they are cited \citep{chapekis2025pew}.

Existing reviews mostly measure accuracy against the auditing outlet's own articles \citep{ebu2025newsintegrity, bbc2025representation} or measure citation retrieval from excerpts \citep{jazwinska2025citation}, while censorship audits often use fixed and historical prompts \citep{naseh2025r1dacted, pan2026censorship}. This project addresses these limitations by examining LLM outputs when queried about recent events and comparing the factual accuracy, depth, and coverage of model-generated information against a fact-checked, multi-outlet baseline. As reliance on LLMs for information grows, these questions have implications for public trust and democratic participation \citep{ebu2025newsintegrity}, particularly given that those most exposed to AI news are also the most likely to trust it \citep{niemanlab2026chatbots}. By studying how information present in trusted online sources may be omitted, reframed, or replaced, it will offer a better understanding of how these channels mediate, frame, or filter information and inform users, developers, and regulators about the reliability of AI-generated news content.

We propose a cross-model audit of how five widely used chatbots present \hl{[*one to three]} recent news events by comparing model outputs against fact-checked reporting from established outlets. Previous work shows that AI assistants frequently misrepresent and mis-cite news \citep{ebu2025newsintegrity, jazwinska2025citation} and that LLM outputs are less epistemically diverse than a basic web search \citep{wright2025epistemic}. However, many studies mainly measure single metrics using archival or adversarial prompts through APIs or proprietary methods. We need a good understanding of how the factualness, sourcing and framing of information \hl{[*check what we are measuring in Methods]} about recent events differ across models and settings to assess whether AI-mediated news narrows or distorts the information environment relative to direct browsing. We will test \hl{[*if search-enabled responses are more accurate than search-disabled ones, if citation support is incomplete, and if low-coverage details of more current events are the most vulnerable to low-quality sources} \citep{newsguard2025monitor, robotstxt2025misinfo}. The results would inform designs that better incorporate uncertainty and source transparency, provide regulators with an empirical foundation for systemic risk assessments and questions of liability for AI-generated claims, and support news literacy guidance on cross-checking AI answers. We will release our code and data.

The remainder of the proposal will discuss relevant literature, experimental and evaluation methods, as well as potential ethical issues.